\documentclass[11pt,a4paper]{article}
\usepackage[margin=0.75in]{geometry}
\usepackage{booktabs}
\usepackage[table]{xcolor}
\definecolor{harpgray}{RGB}{225,225,225}
% diverging scale for the difficulty-gap table: warm = pooled overstates,
% cool = pooled understates, intensity proportional to |gap| (max |gap| = 13.4)
\definecolor{gapPos}{RGB}{214,96,77}
\definecolor{gapNeg}{RGB}{69,117,180}
\usepackage{multirow}
\usepackage{amssymb}
\usepackage{graphicx}
\usepackage{float}
\usepackage{caption}
\title{Session 06 Phase 3 -- Master Results Tables}
\author{}
\date{\today}
\begin{document}
\maketitle

\section*{What each condition means}

Every row below is a different way of turning the model's internal activity while it writes
an answer into a single feature vector. Think of that internal activity as a trail of
numbers -- one per layer, per generated token -- and each condition as a different way of
boiling that trail down.

\begin{description}
\item[core\_max] The baseline. Takes the single highest value each internal signal ever
  reaches while generating the answer -- a snapshot of the loudest moment.
\item[q\_velocity] Instead of the activity itself, tracks how much it \emph{changes} from
  one layer to the next -- speed instead of position. Picks up on whether the model's
  internal representation is still being revised as information moves through the network,
  not just where it ends up.
\item[q\_static] The same kind of peak snapshot as core\_max, but softer: uses a high/low
  percentile instead of the literal maximum/minimum, so one unusual token can't skew the
  whole signal.
\item[core\_concat] Staples core\_max's peak snapshot and q\_velocity's change signal
  together into one longer feature vector -- ``where it ended up'' plus ``how it got
  there.''
\item[joint\_tensor] Instead of computing two signals separately and combining them
  afterward, fits one shared compression across the combined raw data from the start, so
  the signals can influence each other during compression rather than only being stapled
  together at the end.
\item[triple\_concat] Staples all three single signals -- core\_max, q\_static, and
  q\_velocity -- together. The ``use everything we have'' condition.
\end{description}

% ---------- TABLE 1: literature-comparable ----------
\begin{table}[h]
\centering\footnotesize
\begin{tabular}{llcccc}
\toprule
\textbf{Model} & \textbf{Condition} & \textbf{TruthfulQA} & \textbf{TriviaQA} & \textbf{NQ-Open} & \textbf{TyDiQA-GP} \\
\midrule
\multirow{7}{*}{LLaMA-3.1-8B}
  & \cellcolor{harpgray}HARP (published, supervised) & \cellcolor{harpgray}88.5 & \cellcolor{harpgray}92.9 & \cellcolor{harpgray}89.4 & \cellcolor{harpgray}86.6 \\
  & core\_max      & 84.4 & 86.4 & 85.7 & 83.1 \\
  & q\_velocity    & \textbf{86.2} & 85.8 & 84.6 & 82.6 \\
  & q\_static      & 84.3 & 86.7 & 86.1 & \textbf{84.0} \\
  & core\_concat   & 86.0 & 86.8$^{\dagger}$ & 85.3 & 83.2 \\
  & joint\_tensor  & 84.5 & 86.4 & \textbf{86.3} & 83.3 \\
  & triple\_concat & 85.5 & \textbf{87.0} & 85.7 & 83.2 \\
\midrule
\multirow{7}{*}{Qwen2.5-7B}
  & \cellcolor{harpgray}HARP (published, supervised) & \cellcolor{harpgray}88.1 & \cellcolor{harpgray}92.8 & \cellcolor{harpgray}84.0 & \cellcolor{harpgray}88.4 \\
  & core\_max      & 87.1 & 91.3$^{\dagger}$ & 78.8 & \textbf{88.3}$^{\ddagger}$ \\
  & q\_velocity    & 87.7 & 90.9 & 78.2 & 87.0 \\
  & q\_static      & 87.1 & 91.7 & 79.1 & 87.5 \\
  & core\_concat   & 87.8 & 91.5 & \textbf{79.4} & 87.8 \\
  & joint\_tensor  & 87.2 & \textbf{91.8} & 78.5 & 87.0 \\
  & triple\_concat & \textbf{87.8} & 91.7 & 79.3 & 87.9 \\
\bottomrule
\end{tabular}
\caption{\textbf{Literature-comparable results: pooled AUROC (\%) under HARP's evaluation
protocol} (train on 75\% of known questions; test on the rest plus all fully-hallucinated
questions), 5-seed mean, Random Forest readout. All non-gray rows are \emph{our}
unsupervised feature variants scored under this protocol; the gray row is the published
score of HARP's own supervised detector, copied from their paper for reference (not
recomputed by us). Bold = best of our conditions per column (HARP row excluded).
$^{\dagger}$On TriviaQA -- and \emph{only} on TriviaQA -- logistic regression outperforms the
forest, for \emph{both} models, which makes this a property of the dataset rather than the
architecture. LLaMA: core\_max 87.3~$\pm$~0.2, q\_static 87.3~$\pm$~0.3, core\_concat
87.3~$\pm$~0.2. Qwen: q\_static 92.7~$\pm$~0.2, core\_concat 92.6~$\pm$~0.2, joint\_tensor
92.5~$\pm$~0.1, triple\_concat 92.5~$\pm$~0.1, core\_max 92.4~$\pm$~0.1 -- LR wins all six
conditions, and q\_static sits 0.11 points below HARP's published 92.8, our closest approach to
the published reference anywhere in the study. $^{\ddagger}$Qwen/TyDiQA-GP is the
first cell where an unsupervised condition of ours reaches the published supervised score:
core\_max 88.3~$\pm$~1.8 against HARP's reported 88.4, a gap of 0.1 points, well inside our
seed-to-seed spread.}
\end{table}

% ---------- TABLE 2: our internal primary metric ----------
\begin{table}[h]
\centering\small
\begin{tabular}{llcccc}
\toprule
\textbf{Model} & \textbf{Condition} & \textbf{TruthfulQA} & \textbf{TriviaQA} & \textbf{NQ-Open} & \textbf{TyDiQA-GP} \\
\midrule
\multirow{6}{*}{LLaMA-3.1-8B}
  & core\_max      & 81.77 & 81.97 & 76.10 & 84.48 \\
  & q\_velocity    & 83.58 & 81.17 & 75.44 & 85.68 \\
  & q\_static      & 83.28 & 81.97 & 76.72 & 85.26 \\
  & core\_concat   & \textbf{84.35}$^{*}$ & 82.06 & \textbf{76.74} & 85.27 \\
  & joint\_tensor  & 83.76 & 81.84 & 75.70 & 85.13 \\
  & triple\_concat & 83.28 & \textbf{82.27} & 76.01 & \textbf{85.76} \\
\midrule
\multirow{6}{*}{Qwen2.5-7B}
  & core\_max      & 82.47 & 88.56 & 74.45 & 87.39 \\
  & q\_velocity    & 84.33 & 88.11 & 74.88 & 88.50 \\
  & q\_static      & 83.28 & 88.70 & 75.32 & 86.88 \\
  & core\_concat   & \textbf{84.43} & 88.55 & 76.18 & 88.85 \\
  & joint\_tensor  & 82.69 & \textbf{88.80} & 75.87 & 87.11 \\
  & triple\_concat & 84.05 & 88.64 & \textbf{77.61} & \textbf{89.14} \\
\bottomrule
\end{tabular}
\caption{\textbf{Within-prompt AUROC (\%), grouped 5-fold CV, Random Forest.} Among the 10 answers
to the \emph{same} question, does the detector score the hallucinated ones above the truthful
ones? Formally: concordant (hallucinated, truthful) pairs over total such pairs, counting only
pairs drawn from the same question, ties at one half. No HARP row appears because no prior work in
this area reports a per-question metric -- see the preceding section for why we introduce one and
what it costs. Bold = best per column, but note that a bolded lead is
\emph{nominal} unless marked $^{*}$: on TriviaQA, NQ-Open, and TyDiQA-GP no condition's
within-prompt gain over core\_max survives a paired bootstrap (all 95\% CIs include zero;
triple\_concat's TriviaQA lead, for instance, is $\Delta=+0.30$~pts, CI $[-0.14, +0.73]$).
$^{*}$core\_concat vs.\ core\_max on TruthfulQA remains, across all four completed datasets,
the project's \emph{only} statistically significant within-prompt improvement (paired
$\Delta=+2.57$~pts, 95\% CI $[+0.03, +5.26]$, excludes zero). Separately, all four
combination conditions \emph{do} significantly beat q\_velocity on TriviaQA on both metrics.
Because the leads are nominal, the practically useful reading of this table is the \emph{level}
each model reaches, not which condition tops a column.}
\end{table}

% ---------- TABLE 3: the difficulty gap ----------
\begin{table}[h]
\centering\small
\begin{tabular}{llcccc}
\toprule
\textbf{Model} & \textbf{Condition} & \textbf{TruthfulQA} & \textbf{TriviaQA} & \textbf{NQ-Open} & \textbf{TyDiQA-GP} \\
\midrule
\multirow{7}{*}{LLaMA-3.1-8B}
  & core\_max      & \cellcolor{gapPos!25}$+5.15$ & \cellcolor{gapPos!17}$+3.44$ & \cellcolor{gapPos!50}$+10.34$ & \cellcolor{gapNeg!12}$-2.39$ \\
  & q\_velocity    & \cellcolor{gapPos!23}$+4.87$ & \cellcolor{gapPos!17}$+3.53$ & \cellcolor{gapPos!53}$+10.96$ & \cellcolor{gapNeg!17}$-3.52$ \\
  & q\_static      & \cellcolor{gapPos!20}$+4.22$ & \cellcolor{gapPos!18}$+3.67$ & \cellcolor{gapPos!48}$+10.00$ & \cellcolor{gapNeg!11}$-2.31$ \\
  & core\_concat   & \cellcolor{gapPos!20}$+4.22$ & \cellcolor{gapPos!17}$+3.54$ & \cellcolor{gapPos!49}$+10.18$ & \cellcolor{gapNeg!12}$-2.58$ \\
  & joint\_tensor  & \cellcolor{gapPos!21}$+4.42$ & \cellcolor{gapPos!17}$+3.61$ & \cellcolor{gapPos!52}$+10.81$ & \cellcolor{gapNeg!12}$-2.44$ \\
  & triple\_concat & \cellcolor{gapPos!25}$+5.25$ & \cellcolor{gapPos!17}$+3.55$ & \cellcolor{gapPos!52}$+10.72$ & \cellcolor{gapNeg!15}$-3.08$ \\
\cmidrule(lr){2-6}
  & \textbf{mean}  & \cellcolor{gapPos!23}$\mathbf{+4.69}$ & \cellcolor{gapPos!17}$\mathbf{+3.56}$ & \cellcolor{gapPos!51}$\mathbf{+10.50}$ & \cellcolor{gapNeg!13}$\mathbf{-2.72}$ \\
\midrule
\multirow{7}{*}{Qwen2.5-7B}
  & core\_max      & \cellcolor{gapPos!23}$+4.84$ & \cellcolor{gapPos!12}$+2.40$ & \cellcolor{gapPos!63}$+13.02$ & \cellcolor{gapPos!1}$+0.13$ \\
  & q\_velocity    & \cellcolor{gapPos!21}$+4.36$ & \cellcolor{gapPos!13}$+2.65$ & \cellcolor{gapPos!65}$+13.43$ & \cellcolor{gapNeg!2}$-0.42$ \\
  & q\_static      & \cellcolor{gapPos!21}$+4.42$ & \cellcolor{gapPos!12}$+2.58$ & \cellcolor{gapPos!63}$+13.19$ & \cellcolor{gapPos!6}$+1.31$ \\
  & core\_concat   & \cellcolor{gapPos!20}$+4.17$ & \cellcolor{gapPos!13}$+2.61$ & \cellcolor{gapPos!61}$+12.71$ & \cellcolor{gapNeg!2}$-0.45$ \\
  & joint\_tensor  & \cellcolor{gapPos!26}$+5.47$ & \cellcolor{gapPos!13}$+2.62$ & \cellcolor{gapPos!60}$+12.52$ & \cellcolor{gapPos!4}$+0.77$ \\
  & triple\_concat & \cellcolor{gapPos!22}$+4.51$ & \cellcolor{gapPos!13}$+2.67$ & \cellcolor{gapPos!53}$+11.10$ & \cellcolor{gapNeg!4}$-0.74$ \\
\cmidrule(lr){2-6}
  & \textbf{mean}  & \cellcolor{gapPos!22}$\mathbf{+4.63}$ & \cellcolor{gapPos!12}$\mathbf{+2.59}$ & \cellcolor{gapPos!61}$\mathbf{+12.66}$ & \cellcolor{gapPos!1}$\mathbf{+0.10}$ \\
\bottomrule
\end{tabular}
\caption{\textbf{The difficulty gap: pooled minus within-prompt AUROC (percentage points).} Both
terms come from the \emph{same} grouped 5-fold run and the same Random Forest scores, so the
difference isolates exactly one thing -- how much a condition's pooled number depends on being
allowed to compare answers across different questions. Positive means pooled \emph{overstates}
per-response detection; negative means it understates it.
\\[2pt]
\textbf{Colour key.} Shading is a diverging scale on the signed gap: warm
(\cellcolor{gapPos!35}\phantom{xx}) where pooled overstates, cool
(\cellcolor{gapNeg!35}\phantom{xx}) where it understates, with intensity proportional to
magnitude on a common scale across the whole table (deepest shade $=13.4$ points, the largest gap
observed). Near-white cells are gaps close to zero, where the two metrics agree.
\\[2pt]
The striking pattern is how little the gap moves down a column and how much it moves across one.
Within any model$\times$dataset cell the six conditions span about one point; across datasets the
means run from $-2.7$ to $+12.7$. \textbf{The gap is a property of the benchmark, not of the
feature variant.} NQ-Open is the extreme -- at 89--97\% hallucination rates it has very few mixed
questions, so almost all of its pooled comparisons are cross-question and its pooled score is
inflated by 10--13 points. TyDiQA-GP inverts for LLaMA, where pooling mixes questions across
languages and adds variance the within-question comparison never sees; for Qwen the same dataset
sits at essentially zero.
\\[2pt]
Practical consequence: a pooled AUROC cannot be interpreted without knowing its benchmark's gap.
Two published numbers of 86.5 mean different things if one comes from NQ-Open and the other from
TyDiQA-GP.
\\[2pt]
\emph{Note:} Figure~3's right panel reports the same quantity for LLaMA and differs by
0.1--0.5 points ($+4.9$ / $+3.4$ / $+10.8$ / $-2.5$ against this table's core\_max row). Both are
computed from the same features and folds; the small discrepancy is an implementation difference
between the two code paths and is under investigation. This table is built directly from the
Phase-3 result files, consistent with Tables~1 and~2.}
\end{table}

% ---------- TABLE 4: RF vs LR winner grid ----------
\begin{table}[h]
\centering\footnotesize
\begin{tabular}{ll cc cc cc cc}
\toprule
& & \multicolumn{2}{c}{\textbf{TruthfulQA}} & \multicolumn{2}{c}{\textbf{TriviaQA}}
  & \multicolumn{2}{c}{\textbf{NQ-Open}} & \multicolumn{2}{c}{\textbf{TyDiQA-GP}} \\
\cmidrule(lr){3-4}\cmidrule(lr){5-6}\cmidrule(lr){7-8}\cmidrule(lr){9-10}
\textbf{Model} & \textbf{Condition} & RF & LR & RF & LR & RF & LR & RF & LR \\
\midrule
\multirow{6}{*}{LLaMA-3.1-8B}
  & core\_max      & \checkmark &            &            & \checkmark & \checkmark &            & \checkmark &            \\
  & q\_velocity    & \checkmark &            & \checkmark &            & \checkmark &            & \checkmark &            \\
  & q\_static      & \checkmark &            &            & \checkmark & \checkmark &            & \checkmark &            \\
  & core\_concat   & \checkmark &            &            & \checkmark & \checkmark &            & \checkmark &            \\
  & joint\_tensor  & \checkmark &            &            & \checkmark & \checkmark &            & \checkmark &            \\
  & triple\_concat & \checkmark &            &            & \checkmark & \checkmark &            & \checkmark &            \\
\midrule
\multirow{6}{*}{Qwen2.5-7B}
  & core\_max      & \checkmark &  &  & \checkmark & \checkmark &  & \checkmark &  \\
  & q\_velocity    & \checkmark &  &  & \checkmark & \checkmark &  & \checkmark &  \\
  & q\_static      & \checkmark &  &  & \checkmark & \checkmark &  & \checkmark &  \\
  & core\_concat   & \checkmark &  &  & \checkmark & \checkmark &  & \checkmark &  \\
  & joint\_tensor  & \checkmark &  &  & \checkmark & \checkmark &  & \checkmark &  \\
  & triple\_concat & \checkmark &  &  & \checkmark & \checkmark &  & \checkmark &  \\
\bottomrule
\end{tabular}
\caption{\textbf{Which classifier wins, under the HARP protocol} -- \checkmark marks the
higher of RF/LR (5-seed mean pooled AUROC) for that condition$\times$dataset cell. RF sweeps
TruthfulQA, NQ-Open, and TyDiQA-GP (18/18 LLaMA cells). TriviaQA flips: LR wins 5 of 6
conditions, RF holds only on q\_velocity. See Table~1's caption for the actual TriviaQA
numbers (LR's wins there are narrow, 0.2--0.9pt -- not a landslide like RF's wins
elsewhere). Qwen reproduces this exactly: RF sweeps TruthfulQA, NQ-Open and TyDiQA-GP (18/18
cells) and LR sweeps all six TriviaQA conditions. Across both models, 42 of 48
condition$\times$dataset cells favour RF and all 12 exceptions are TriviaQA. The inversion
therefore tracks the \emph{dataset}, not the architecture.}
\end{table}

\clearpage
\section*{Within-prompt AUROC: motivation, definition, and why we introduce it}

\subsection*{The convention, and what it silently assumes}

Hallucination detection is evaluated almost universally with a single dataset-level AUROC. HARP
states that ``AUROC (area under the ROC curve) is employed as the primary evaluation metric'';
INSIDE that ``AUROC is a popular metric to evaluate the quality of a binary classifier''; and
Janiak et al.\ that AUROC ``assesses the ability of a hallucination detection method to correctly
rank positive and negative instances (hallucinations vs.\ non-hallucinations).''

All three generate \emph{multiple answers per question} -- ten, in every case, as do we. So the
data has an explicit two-level structure: answers nested inside questions. Yet none of the three
states how scores from \emph{different} questions are combined, and the standard implementation
simply flattens the array. Every hallucinated answer is ranked against every truthful answer in
the dataset, regardless of whether they answer the same question. The grouping is present in the
data and discarded without comment.

\subsection*{Why that matters: the shortcut}

Questions differ enormously in difficulty. In our data, 300 of TruthfulQA's 817 questions are
ones where the model got all ten attempts wrong, and 111 where it got all ten right. Under pooled
AUROC, every answer from the first group is compared against every answer from the second, and a
detector gets all of those pairs correct simply by outputting a higher score for hard questions.
It never has to distinguish a good answer from a bad one.

This is not hypothetical. Janiak et al.\ report two ``simple baselines'' -- \texttt{Mean-Len}
(mean answer length across a question's ten generations) and \texttt{Std-Len} (their standard
deviation) -- that are \emph{per-question constants}: identical for all ten answers to a question,
carrying literally zero information about any individual answer. On LLaMA they average 0.721 and
0.751 AUROC respectively, against 0.747 for Eigenscore, a sophisticated internal-state detector.
A feature that \emph{cannot} discriminate between two answers to the same question, even in
principle, matches or beats the state of the art under pooled AUROC. They present this as evidence
that simple baselines are competitive; we read it as evidence that the metric is measuring
something other than what it appears to.

\subsection*{Definition}

Both metrics are computed from the \emph{same} score vector -- one detector, trained once, scoring
every answer out-of-fold. They differ only in which comparisons are permitted.

AUROC in either form is a pair-counting quantity: among all (hallucinated, truthful) pairs, the
fraction in which the hallucinated answer received the higher score, with ties counting one half.

\begin{itemize}
\item \textbf{Pooled AUROC} forms every such pair in the dataset, including pairs drawn from
different questions.
\item \textbf{Within-prompt AUROC} forms only pairs whose two members answer the \emph{same}
question, then aggregates concordant pairs over total pairs across all questions.
\end{itemize}

\noindent A question whose ten answers are all truthful, or all hallucinated, yields no valid pair
and drops out entirely. Only \emph{mixed} questions -- where the model both succeeded and failed
on the same input -- contribute.

\textbf{``OOF'' means out-of-fold.} Under \texttt{GroupKFold(5)} grouped by prompt, all ten answers
to a question stay together, so a question is never split across training and test. Each answer is
scored by the one model that did not see it. ``Pooled OOF AUROC'' is a single AUROC over all
held-out scores, not an average of five per-fold values.

\subsection*{Why this is the right correction, and not a novel one}

Conditioning on a grouping variable to remove a confound is standard practice elsewhere.
Information retrieval computes MAP and NDCG per query and then averages; pooling documents across
queries would be considered plainly wrong. Recommender systems use per-user AUC (``GAUC'') for
exactly the reason we need it here: pooled AUC conflated ``this user clicks often'' with ``this
item suits this user.'' Substituting question for user and answer for item gives our setting
exactly.

We therefore do not claim within-prompt AUROC as a new metric. We claim that hallucination
detection has not been applying it, that the two-level structure required for it is already
present in every standard benchmark, and that the difference it makes is large -- which is what
Figure~3 quantifies.

\subsection*{What it costs}

The correction is not free, and the limitations should be stated plainly.

\begin{itemize}
\item \textbf{Homogeneous questions are discarded.} On Qwen/NQ-Open only 348 of 3{,}610 questions
are mixed, so the metric rests on 5{,}365 pairs rather than the full dataset, and its confidence
intervals are correspondingly wide. Where the model almost never succeeds, there is little to
measure.
\item \textbf{It is not comparable to published numbers.} No prior work reports it, so pooled
AUROC must still be reported alongside for literature comparison. We report both throughout.
\item \textbf{It is not uniformly stricter.} On TyDiQA-GP within-prompt AUROC \emph{exceeds}
pooled. Pooling there mixes questions across languages, adding cross-question variance that the
within-question comparison never sees. The sign of the gap is itself a per-benchmark diagnostic,
not a fixed penalty.
\end{itemize}

\noindent The gap between the two numbers is the quantity of interest, and Figure~3 decomposes it
directly.

\clearpage
\section*{Figure 1 -- where in the network the signal lives}

\begin{figure}[H]
\centering
\includegraphics[width=\linewidth]{{fig1_layer_divergence_llama-3.1-8b_all}.pdf}
\caption*{\small LLaMA-3.1-8B. Left: how much larger the average size of the internal activity is
for hallucinated answers than for truthful ones, in pooled standard deviations. Right: detection
score when the detector is allowed to use only that one layer (64 components, logistic
regression). Bands are 95\% confidence intervals. Only layers 15--23 exist on disk.}
\end{figure}

\subsection*{What it shows}
\begin{itemize}
\item \textbf{Left panel: the curves are flat.} There is no bump at any layer. The ``layer-22
hump'' we have been assuming does not exist -- the whole range of movement across nine layers is
about 0.1 standard deviations, which is nothing.

\item \textbf{The datasets sit far apart, and one has the opposite sign.} NQ-Open is at $+1.3$,
TriviaQA $+0.7$, TyDiQA-GP $+0.35$, and TruthfulQA at $-0.85$ -- there, hallucinated answers have
\emph{smaller} internal activity, not larger. So this quantity behaves like a fingerprint of the
dataset, not a signal of hallucination.

\item \textbf{Critically: this quantity does not predict detection at all.} NQ-Open has the
largest gap and TruthfulQA has a negative one, yet both sit at the top of the right panel. Since
the 15--23 window was originally chosen using exactly this kind of energy audit, the method we
used to pick our window does not track what we actually care about.

\item \textbf{Right panel: layer 15 is the best single layer on all four datasets.} Three of the
four decline steadily to the right; TruthfulQA is flat rather than falling.

\item \textbf{Layer 15 is the left edge of our window, which is a warning sign.} When the best
value sits on the boundary of the range you looked at, the true best is usually outside it. The
honest reading is that we may have been extracting from the wrong part of the network, and we
cannot tell because the pipeline crops at 15. Settling this needs one GPU pass over earlier
layers.

\item \textbf{A single layer appears to match the entire pipeline.} Layer 15 alone, with 64
numbers and a simple linear model, scores about 0.877 / 0.862 / 0.873 / 0.830 across the four
datasets, against 0.869 / 0.854 / 0.865 / 0.821 for the full 9-layer, 320-number random forest.
The confidence intervals overlap, so the claim is ``matches,'' not ``beats'' -- but even a match
means the other eight layers are contributing close to nothing. This needs one cheap paired check
before we state it publicly.
\end{itemize}

\clearpage
\section*{Figure 2 -- do the static and velocity streams carry different things?}

\begin{figure}[H]
\centering
\includegraphics[width=\linewidth]{{fig2_principal_angles_llama-3.1-8b_all}.pdf}
\caption*{\small Left: for each pair of directions (1st most important, 2nd, and so on up to
64th), how closely aligned the static and velocity subspaces are. 1.0 means the two directions are
identical; 0 means completely unrelated. Right: how stable the static subspace is, comparing folds
within a dataset (diagonal) and datasets against each other (off-diagonal).}
\end{figure}

\subsection*{What it shows}
\begin{itemize}
\item \textbf{Three clearly separated bands in the left panel.} The black line (the two halves of
the \emph{same} stream, q95 vs q05) stays near 1.0 almost the whole way -- they are near-copies,
sharing about 58 of 64 directions. The coloured lines (static vs velocity) drop below 0.5 around
the 15th direction -- about 15 of 64 shared. The grey line (two random subspaces) never rises
above 0.18 -- nothing shared.

\item \textbf{So static and velocity really are geometrically different} -- far more different
than chance, far less different than two halves of one stream. That part of the story holds.

\item \textbf{But being different did not make them useful together.} Combining the two streams
buys only 0.1--2.5 points, mostly not statistically significant. So the two streams point in
different directions while carrying largely the \emph{same} information about hallucination. This
is the honest headline, and it explains why several rounds of feature engineering produced about
one point.

\item \textbf{Where the overlap sits matters.} The first five directions are nearly identical
(above 0.8). The part the two streams share is their biggest, most dominant part -- which is
plausibly the part a classifier leans on most. The genuinely different directions are the small,
tail ones.

\item \textbf{A free simplification.} Because q05 is nearly a copy of q95, we could likely drop
q05 entirely and halve the feature size at no cost. That is a CPU-only check on data we already
have.

\item \textbf{Right panel, and a good rigour point:} comparing two different folds of the same
dataset gives 0.99--1.00. Our unsupervised basis is essentially deterministic -- it does not wobble
when the data is resampled. Across different datasets it is 0.69--0.92, so there is substantial
shared structure, with TyDiQA-GP the most distinct.

\item \textbf{One caution when presenting:} the diagonal and off-diagonal cells of that heatmap
measure two different comparisons. It is stated in the title but is easy to miss.
\end{itemize}

\clearpage
\section*{Figure 3 -- what the standard metric is actually measuring}

\begin{figure}[H]
\centering
\includegraphics[width=\linewidth]{{fig3_difficulty_decomposition_llama-3.1-8b_all}.pdf}
\caption*{\small Left: pooled AUROC using full features, using only each question's 10-answer
average (``centroid-only,'' which destroys all per-answer information), and using each answer
minus that average (``delta-only,'' which destroys all question-identity information). Right:
pooled AUROC minus within-prompt AUROC, per dataset.}
\end{figure}

\subsection*{What it shows}
\begin{itemize}
\item \textbf{Question identity alone gets most of the way there.} The centroid-only bars reach
0.71--0.81. Those features contain \emph{zero} information about any individual answer -- every
answer to a question has been replaced by the same average. Yet they recover 66--85\% of the
pooled score's lift above chance.

\item \textbf{What this does and does not mean.} It does not mean the metric is invalid -- those
features still come from the model's own internal states. It means pooled AUROC is mostly
measuring \emph{which questions the model will fail on}, not \emph{which specific answers are
lies}. Most readers of a hallucination-detection paper assume the second.

\item \textbf{The size of the distortion varies a lot by benchmark} (right panel). NQ-Open is the
most inflated at $+0.108$, then TruthfulQA $+0.049$ and TriviaQA $+0.034$. TyDiQA-GP is the only
one where pooled \emph{understates} performance, at $-0.025$.

\item \textbf{The most useful result is not currently drawn in the figure.} Removing the question
average and keeping only the per-answer part \emph{beats} the full features on within-prompt
AUROC: 83.2 vs 82.0 on TruthfulQA, 83.8 vs 82.0 on TriviaQA. So the question average is not just
unhelpful for per-answer detection -- it is actively getting in the way.

\item \textbf{That is a plain subtraction beating all six engineered feature variants} on the
metric we care about, at no compute cost. It should go into the condition grid for both models.

\item \textbf{One honest caveat on it.} Subtracting the average needs the other nine answers
available at test time, so it is a ``multi-sample'' method -- the same setting as SelfCheckGPT and
INSIDE, which is legitimate and precedented, but it is not a like-for-like comparison against
single-pass HARP and must be labelled as such. It uses no labels, so there is no leakage.

\item \textbf{This lines up with the tables.} Combination conditions beat the baseline on pooled
but never on within-prompt, for both models and every dataset -- exactly what you would expect if
those combinations are adding question-difficulty information rather than per-answer signal. Two
separate analyses reaching the same conclusion.
\end{itemize}

\clearpage
\section*{Novel contributions}

Claims are graded by how well the current data supports them. \textbf{[Established]} means our own
measurements support it directly. \textbf{[Needs lit check]} means the result is solid but we have
not confirmed no one has published it. \textbf{[Do not claim]} records things we deliberately are
\emph{not} asserting, so the boundary is explicit.

\subsection*{Part 1 -- What we have now}

\paragraph{1. Pooled AUROC is mostly a question-difficulty detector, and we measured how mostly.}
\textbf{[Established, needs lit check]} Replacing every answer by its question's 10-answer mean
destroys all per-answer information, yet those features recover \textbf{66--85\% of pooled AUROC's
lift above chance} across four datasets. The centroid/delta decomposition is the instrument, and
it is the contribution -- not the observation that a confound exists, but a number for how large
it is. Janiak et al.\ independently produced supporting evidence without framing it this way:
their \texttt{Std-Len} baseline, a per-question constant that cannot separate two answers to the
same question even in principle, scores 0.751 pooled AUROC against Eigenscore's 0.747.

\paragraph{2. The difficulty gap is a property of the benchmark, not of the method.}
\textbf{[Established]} Across 48 model$\times$condition$\times$dataset cells, the pooled-minus-
within-prompt gap varies by about one point across the six conditions within a cell, and by
\emph{fifteen} points across datasets ($-2.7$ to $+12.7$; Table~3). The practical consequence is
concrete and checkable: two published pooled AUROCs of 86.5 describe different things if one comes
from NQ-Open and the other from TyDiQA-GP. This motivates reporting the gap as a benchmark
descriptor -- see Part 2.

\paragraph{3. The sign of the gap varies, so this is a diagnostic and not a correction factor.}
\textbf{[Established]} TyDiQA-GP inverts for LLaMA: pooled \emph{understates} per-response
detection there, plausibly because pooling mixes questions across languages. Any framing that
treats pooled AUROC as uniformly inflated is wrong.

\paragraph{4. Energy divergence does not track discriminability, which invalidates a common
layer-selection heuristic.} \textbf{[Established, needs lit check]} Figure~1: the norm gap between
hallucinated and truthful answers is essentially flat across the extraction window
($\sim$0.1~SD over nine layers), varies enormously between datasets ($+1.3$ to $-0.85$),
\emph{flips sign} on TruthfulQA, and is uncorrelated with per-layer detection AUROC -- NQ-Open has
the largest norm gap and TruthfulQA a negative one, yet both rank highest on discriminability. Our
own extraction window was selected by exactly such an audit, so this is a self-critique with
general force.

\paragraph{5. Geometric complementarity does not imply informational complementarity.}
\textbf{[Established]} Figure~2: the static and velocity subspaces share only $\sim$16 of 64
directions against a random null of 0 and a within-stream reference of $\sim$58 -- genuinely
different geometry -- yet combining them buys 0.1--2.5 points, mostly not significant. The overlap
concentrates in the leading, highest-variance directions, which is plausibly what a classifier
leans on. This explains why several rounds of feature engineering produced about one point, and it
is a more useful negative result than ``the combination did not help.''

\paragraph{6. A plain mean-subtraction outperforms all six engineered feature variants.}
\textbf{[Established on 2 of 4 datasets; extension pending]} Removing each answer's prompt centroid
beats the full features on within-prompt AUROC (83.2 vs 82.0 on TruthfulQA, 83.8 vs 82.0 on
TriviaQA) at zero compute cost. Must be reported as a \emph{multi-sample} method, since centering
requires the sibling answers at inference -- the SelfCheckGPT/INSIDE setting, legitimate and
precedented, but not comparable to single-pass HARP.

\paragraph{7. An unsupervised pipeline reaches published supervised performance.}
\textbf{[Established]} On Qwen2.5-7B, three of four datasets sit within 0.3 points of HARP's
published supervised scores: TyDiQA-GP $-0.1$, TriviaQA $-0.11$, TruthfulQA $-0.3$. The best
condition on two of them is plain \texttt{core\_max}, which supports the reading that the
\emph{pipeline} is the contribution rather than any particular feature variant.

\paragraph{8. Two smaller consistency results.} \textbf{[Established]} The RF/LR inversion tracks
the dataset, not the architecture: 42 of 48 cells favour Random Forest and all 12 exceptions are
TriviaQA, for both models. And condition ranking is stable across architectures with different
hidden dimensions (4096 vs 3584) -- \texttt{core\_concat}/\texttt{triple\_concat} occupy the top two
positions in 5 of 6 completed cells.

\paragraph{What we explicitly do not claim.} \textbf{[Do not claim]} Within-prompt AUROC is
\emph{not} a new metric -- information retrieval computes MAP and NDCG per query, and recommender
systems use per-user AUC (GAUC) for precisely this confound. Our claim is that hallucination
detection has not been applying it despite every standard benchmark already having the required
structure. Nor do we claim to have discovered that hallucination benchmarks are confounded;
Janiak et al.\ established that for labels and response length. Ours is a different axis.

\subsection*{Part 2 -- Future contributions reachable from the current state}

Ordered by expected value per unit cost. All but the last two are CPU-only on artifacts already on
disk.

\paragraph{A. Re-label with an LLM judge and re-run everything. [GPU-free, API cost. Highest
value.]} Janiak et al.\ show ROUGE-based correctness labels have 0.401 precision against human
judgment, and that detector rankings change substantially -- up to 45.9\% -- when labels are
corrected. \emph{Our labels are ROUGE-L $>0.7$ or BLEURT $>0.5$, so we inherit this criticism
in full.} Every generation is already pinned to disk and labelling is a separate pipeline stage,
so re-labelling and re-running costs no GPU time. This is simultaneously our largest exposure and
the cheapest way to close it: if our findings survive human-aligned labels, they become far more
defensible, and ``we re-verified under LLM-as-Judge'' is itself a contribution.

\paragraph{B. Does one layer match the whole pipeline? [CPU, hours.]} Figure~1(ii) suggests layer
15 alone, at 64 dimensions with logistic regression, matches the full 9-layer 320-dimension random
forest on all four datasets (0.877/0.862/0.873/0.830 vs 0.869/0.854/0.865/0.821). Confidence
intervals overlap, so the verifiable claim is ``matches,'' not ``beats'' -- but a match means eight
of nine layers contribute nothing and the compression story strengthens fivefold. Requires one
paired test against \texttt{core\_max} on identical folds.

\paragraph{C. Is the extraction window in the wrong place? [One GPU pass. Highest ceiling.]} The
best single layer is 15 -- the \emph{left edge} of the window. An optimum sitting on a boundary
usually means the true optimum lies outside it. If the signal peaks before layer 15, then the
convention of extracting from middle-to-late layers may be misplaced across this whole line of
work, and every number we have is a lower bound. One extraction pass over layers $\sim$7--15 on two
datasets settles it. Note the caveat: single-layer AUROC is not marginal contribution inside a
9-layer fit, so a monotone profile is suggestive rather than conclusive.

\paragraph{D. Promote delta-centering from diagnostic to method. [CPU, free.]} Extend
contribution~6 across the full grid and both models. If it holds, the paper gains a concrete
improvement that follows directly from the decomposition -- the diagnostic and the method share one
idea, which is a much tighter story than reporting them separately.

\paragraph{E. Report the gap as a benchmark descriptor. [CPU, near-free.]} Turn contribution~2
into a proposed practice: alongside any pooled AUROC, report that benchmark's difficulty gap so
readers can interpret the number. We would publish gaps for the four standard benchmarks as a
reference table. This is the cheapest path from ``we found a problem'' to ``we gave the community
a fix.''

\paragraph{F. Aggregate properly instead of testing cell by cell. [CPU, hours.]} Combination
conditions beat the baseline on pooled but never on within-prompt in \emph{any single cell}. A
random-effects meta-analysis across all eight cells, plus a Friedman/Nemenyi test over the six
conditions, resolves whether that is a real null or eight underpowered tests. Either outcome is
publishable, and ``this subfield tests underpowered comparisons one dataset at a time'' is a
methods point in its own right.

\paragraph{G. Halve the feature dimensionality for free. [CPU, hours.]} Figure~2 shows the q95 and
q05 halves share $\sim$58 of 64 directions. If dropping q05 costs nothing measurable, every feature
variant halves in size. Minor alone; useful combined with B.

\paragraph{H. Cross-dataset transfer. [CPU, heavy.]} Fit on one dataset, evaluate on another, all
16 pairs per model, against HARP's published transfer matrix. Tests whether the learned subspace is
general or dataset-specific -- and Figure~2(B) already gives a geometric prediction to check it
against (cross-dataset subspace similarity 0.69--0.92, TyDiQA-GP most distinct), so the two
analyses can corroborate each other.

\paragraph{Sequencing.} A is the priority: until labels are human-aligned, every other contribution
carries an asterisk. B, D, E and F are cheap and can run in parallel with it. C has the highest
ceiling but the highest cost, and is best attempted once A has confirmed the findings are real.

\end{document}
